\section{从 Fourier 级数到连续变换}
\label{sec:08-Numerical-Analysis/08-Fourier-and-Spectral-Methods/01}

\subsection{时域里的困境：你看到的和你想找的不是同一层}

一个周期信号——一天的温度曲线、一段语音、一台电机的振动读数——同时裹着缓慢的大趋势和快速的小波动。时域里每个函数值把不同频率的贡献全搅在一起。你盯着波形图，能感觉到``这里有个大周期''和``上面叠着毛毛躁躁的小波动''，但仅凭肉眼无法把它们定量拆开。

更麻烦的是两种最常见的运算：

\begin{itemize}
\tightlist
\item
  \textbf{微分}：\(\frac{d}{dt}\sin(10t) = 10\cos(10t)\)，高频被放大一个数量级。如果你的信号里混着测量噪声（富含高频），直接做数值微分就是在放大噪声。
\item
  \textbf{卷积}：\(t\) 时刻的卷积值 \(\int f(\tau)g(t-\tau)d\tau\) 把整个时间轴的值都卷进来，所有时刻彼此耦合，没法拆成``这个频率过这个滤波器、那个频率过那个滤波器''。
\end{itemize}

滑动平均能粗略地把快慢分开——窗口越宽，剩下越慢的趋势分量。但窗宽是人为选的，20 个点的窗口和 50 个点的窗口给出的是两套快慢划分，没有天然的正确答案。

你想要的不只是``一次平滑''，而是一套坐标变换：把函数从时域旋转到一组能独立处理各频率分量的坐标轴上，做完微分和滤波再转回来。

\subsection{复指数：微分算子的特征函数}

看一个线索。对正弦求导：\(\frac{d}{dx}\sin(kx) = k\cos(kx)\)。频率 \(k\) 变成了振幅的倍数——不同的频率在微分下有不同的响应。这暗示：如果能把函数写成不同频率成分的叠加，微分就只是对每个成分分别乘以一个倍数。

复指数 \(e^{ikx}\) 把这件事做到极致。它是微分算子 \(\frac{d}{dx}\) 的特征函数，特征值 \(ik\)。求导不改变它的形状，只改变它的振幅和相位。周期 \(2\pi\) 的区间上，不同频率的复指数彼此正交：

\[
\frac{1}{2\pi}\int_{-\pi}^{\pi}e^{ikx}e^{-i\ell x}\,dx = \delta_{k\ell}.
\]

这意味它们构成 \(L^2[-\pi,\pi]\) 的一组完备正交基。任何（性质够好的）周期函数都可以在这组基上展开：

\[
f(x) \sim \sum_{k\in\Z}c_k e^{ikx},
\qquad
c_k = \frac{1}{2\pi}\int_{-\pi}^{\pi}f(x)e^{-ikx}\,dx.
\]

低频 \(|k|\) 小，描述缓变趋势；高频 \(|k|\) 大，描述快变细节。这里的``频率''是基向量的方向标签，不是函数之外另加的物理量。实数形式的 \(\sin(kx)\) 和 \(\cos(kx)\) 只是把正负频率配对——\(c_k e^{ikx}+c_{-k}e^{-ikx}\) 重排成 \(a_k\cos(kx)+b_k\sin(kx)\)。

重要的是 Parseval 恒等式：

\[
\frac{1}{2\pi}\int_{-\pi}^{\pi}|f(x)|^2\,dx = \sum_{k\in\Z}|c_k|^2.
\]

时域总能量等于各频率系数平方和。Fourier 展开是 \(L^2\) 空间中的一次正交坐标变换——像旋转向量坐标系，不改变长度。信息没有丢失，只是换了一组更便利的基矢量。

\subsection{光滑性就是高频系数坠崖的速度}

如果 \(f\) 周期可微且边界项匹配（\(f(\pi)=f(-\pi)\)），分部积分一次：

\[
c_k = \frac{1}{ik}\cdot\frac{1}{2\pi}\int_{-\pi}^{\pi}f'(x)e^{-ikx}\,dx.
\]

每多一阶可积的导数，系数前面多挂一个因子 \(1/|k|\)。

具体地：
\begin{itemize}
\tightlist
\item \(f\) 一阶可微：\(|c_k|\sim 1/|k|\)；
\item \(f\) 二阶可微：\(|c_k|\sim 1/k^2\)；
\item \(f\) 无穷可微：\(|c_k|\) 比任何多项式的倒数衰减都快；
\item \(f\) 能解析延拓到复平面中的条带 \(|\operatorname{Im}z|<\alpha\) 内：\(|c_k|\le Ce^{-\alpha|k|}\)，指数衰减。
\end{itemize}

衰减率 \(\alpha\) 恰好是条带的半宽。函数在复平面上能``撑开''多远，频域系数就塌得多快。

这个对应关系是谱方法的基石。一个 \(C^\infty\) 函数的 Fourier 系数在几十项后就可以截断，逼近精度达到机器零——因为漏掉的高频系数总能量可以忽略。反过来，如果你算出一组 Fourier 系数，看到尾部迟迟不掉，要么函数本身有尖角或跳跃（导数不连续），要么周期拼接处有间断（\(f(-\pi)\neq f(\pi)\) 导致隐式跳跃），要么数据里混进了高频噪声。从系数衰减速度诊断信号的``干净程度''，是频域分析最常用也最直白的技巧。

\subsection{收敛有三种意思，混了就会误判}

\(L^2\) 函数：部分和 \(S_N f(x)=\sum_{|k|\le N}c_k e^{ikx}\) 在均方意义下收敛于 \(f\)。但均方收敛不担保任何一点上的函数值是对的。

分段光滑函数在跳跃点处，部分和趋近于左右极限的平均值：

\[
\lim_{N\to\infty}S_N f(x) = \frac{f(x^-)+f(x^+)}{2}.
\]

跳跃点附近出现 Gibbs 振荡：\(N\) 增大时振荡区域越来越窄，但最大过冲的比例不会趋于零——始终维持跳跃幅度的约 \(9\%\)。取再多项都没法令跳点邻域像光滑函数那样干净。

方波是 Gibbs 的典型感染者：系数 \(|c_k|\sim 1/k\)（一阶可微的边界条件，恰好对应跳跃带来的 \(1/k\) 尾部），过冲客观存在，不是画图精度问题。Gibbs 不是 Fourier 级数的``缺陷''——它是用连续基函数的有限线性组合去逼近间断的必然代价，换成任何其他连续基也一样。

三种收敛：\(L^2\) 收敛（能量意义）、逐点收敛（几乎处处成立但对 \(L^2\) 不保证处处成立）、一致收敛（要求系数绝对可和，\(\sum|c_k|<\infty\)，对应较强条件）。在不同应用场景中需要不同级别的收敛保证，混用是频域错误判断的主要来源。

\subsection{从级数到积分：让周期走远}

非周期信号不能直接用级数展开。把观察周期 \(L\) 不断拉大：频率间隔 \(\Delta\omega = 2\pi/L \to 0\)。离散频率和过渡为连续频率积分：

\[
\hat f(\omega) = \int_{-\infty}^{\infty}f(t)e^{-i\omega t}\,dt,
\qquad
f(t) = \frac{1}{2\pi}\int_{-\infty}^{\infty}\hat f(\omega)e^{i\omega t}\,d\omega.
\]

在适当条件下两式互为逆变换。\(2\pi\) 放在哪一侧是约定问题：若用 \(\nu=\omega/2\pi\)（cycles per second），指数写成 \(e^{-2\pi i\nu t}\)，归一化常数会移动到不同的位置。混用 \(\omega\) 和 \(\nu\) 导致 \(2\pi\) 多一个或少一个，是 Fourier 公式使用中最常见的符号错误。任何公式引用前都应先确认采用的归一化约定。

\subsection{频域：算子变成乘法，耦合变成独立}

若 \(f\) 在无穷远处衰减够快（导数变换时边界项消失）：

\[
\widehat{f'}(\omega) = i\omega\hat f(\omega).
\]

微分——一个局部的、需要取极限才能定义的运算——在频域变成逐频率乘以 \(i\omega\)。高频乘以更大的 \(|\omega|\)：所谓``微分放大噪声''现在有了精确的定量表述——噪声谱在高频段被线性放大。

卷积同样变得简单：

\[
(f*g)(t)=\int f(\tau)g(t-\tau)\,d\tau
\;\Longrightarrow\;
\widehat{f*g}(\omega)=\hat f(\omega)\,\hat g(\omega).
\]

时域中每点都需访问整段信号的全局操作，降解为逐频率的独立乘法。反过来说，时域的乘积对应频域的卷积：时间越紧窄（信号能量集中在短时间内），频谱越宽散。Gaussian 函数的 Fourier 变换仍是 Gaussian，且标准差的乘积为常数——紧窄与宽散互为倒数，这是时频不确定性的一种独立于量子力学而存在的数学事实。

这两条性质——微分变乘法、卷积变乘法——是``为什么要把问题搬到频率空间''的最硬的答案。在一组基底下纠缠的耦合关系，在另一组基底下不一定纠缠。Fourier 分析的价值不在于公式本身，而在于它为多数信号处理问题选择了正确的坐标系。

\subsection{普通积分画不出的变换}

常数 \(f(t)=1\) 和纯正弦 \(f(t)=e^{i\omega_0 t}\) 不属于 \(L^1(\R)\)。它们的绝对值积分发散，不能用上面的积分公式直接定义 Fourier 变换。这类函数的变换需要用分布语言描述：常数的变换是 Dirac \(\delta(\omega)\)，纯正弦的变换是频移的 \(\delta\)。\(\delta\) 不是一个普通函数，它在 \(\omega\neq0\) 处``取值''为 0，在 \(\omega=0\) 处``取值''无穷大，但积分等于 1。

一般 \(L^2\) 函数平方可积却不一定绝对可积，同样不能用逐点积分定义变换。Plancherel 定理通过延拓来解决：先在稠密好子空间（Schwartz 函数）上定义 Fourier 变换，再按 \(L^2\) 范数的连续性延拓到全空间。延拓保证变换存在且保持内积，但不能逐点按积分公式理解。

落到实际数值计算，信号永远是有限时窗里采到的离散样本——处理的是离散 Fourier 变换（DFT）和采样定理的世界，不是连续 Fourier 变换。窗截断→频谱泄漏，离散采样→混叠，有限分辨率→栅栏效应。从连续 Fourier 积分走到实际频谱分析的每一步都会掉信息，后续各节逐一拆开这些问题：\hyperref[sec:08-Numerical-Analysis/08-Fourier-and-Spectral-Methods/02]{``采样、混叠与可恢复性''}讨论连续到离散丢失了什么，\hyperref[sec:08-Numerical-Analysis/08-Fourier-and-Spectral-Methods/03]{``DFT 是变换，FFT 是算法''}厘清哪些信息已被不可逆地丢弃。
